% Options for packages loaded elsewhere
\PassOptionsToPackage{unicode}{hyperref}
\PassOptionsToPackage{hyphens}{url}
\documentclass[
]{article}
\usepackage{xcolor}
\usepackage{amsmath,amssymb}
\setcounter{secnumdepth}{-\maxdimen} % remove section numbering
\usepackage{iftex}
\ifPDFTeX
  \usepackage[T1]{fontenc}
  \usepackage[utf8]{inputenc}
  \usepackage{textcomp} % provide euro and other symbols
\else % if luatex or xetex
  \usepackage{unicode-math} % this also loads fontspec
  \defaultfontfeatures{Scale=MatchLowercase}
  \defaultfontfeatures[\rmfamily]{Ligatures=TeX,Scale=1}
\fi
\usepackage{lmodern}
\ifPDFTeX\else
  % xetex/luatex font selection
\fi
% Use upquote if available, for straight quotes in verbatim environments
\IfFileExists{upquote.sty}{\usepackage{upquote}}{}
\IfFileExists{microtype.sty}{% use microtype if available
  \usepackage[]{microtype}
  \UseMicrotypeSet[protrusion]{basicmath} % disable protrusion for tt fonts
}{}
\makeatletter
\@ifundefined{KOMAClassName}{% if non-KOMA class
  \IfFileExists{parskip.sty}{%
    \usepackage{parskip}
  }{% else
    \setlength{\parindent}{0pt}
    \setlength{\parskip}{6pt plus 2pt minus 1pt}}
}{% if KOMA class
  \KOMAoptions{parskip=half}}
\makeatother
\usepackage{longtable,booktabs,array}
\newcounter{none} % for unnumbered tables
\usepackage{calc} % for calculating minipage widths
% Correct order of tables after \paragraph or \subparagraph
\usepackage{etoolbox}
\makeatletter
\patchcmd\longtable{\par}{\if@noskipsec\mbox{}\fi\par}{}{}
\makeatother
% Allow footnotes in longtable head/foot
\IfFileExists{footnotehyper.sty}{\usepackage{footnotehyper}}{\usepackage{footnote}}
\makesavenoteenv{longtable}
\usepackage{graphicx}
\makeatletter
\newsavebox\pandoc@box
\newcommand*\pandocbounded[1]{% scales image to fit in text height/width
  \sbox\pandoc@box{#1}%
  \Gscale@div\@tempa{\textheight}{\dimexpr\ht\pandoc@box+\dp\pandoc@box\relax}%
  \Gscale@div\@tempb{\linewidth}{\wd\pandoc@box}%
  \ifdim\@tempb\p@<\@tempa\p@\let\@tempa\@tempb\fi% select the smaller of both
  \ifdim\@tempa\p@<\p@\scalebox{\@tempa}{\usebox\pandoc@box}%
  \else\usebox{\pandoc@box}%
  \fi%
}
% Set default figure placement to htbp
\def\fps@figure{htbp}
\makeatother
\setlength{\emergencystretch}{3em} % prevent overfull lines
\providecommand{\tightlist}{%
  \setlength{\itemsep}{0pt}\setlength{\parskip}{0pt}}
\usepackage{bookmark}
\IfFileExists{xurl.sty}{\usepackage{xurl}}{} % add URL line breaks if available
\urlstyle{same}
\hypersetup{
  hidelinks,
  pdfcreator={LaTeX via pandoc}}

\author{}
\date{}

\begin{document}

\section{Group Final Report}\label{group-final-report}

\subsection{Accelerating Gemma-3-12B Inference with Standard Speculative
Decoding and
EAGLE-3}\label{accelerating-gemma-3-12b-inference-with-standard-speculative-decoding-and-eagle-3}

\subsection{1. Introduction}\label{1-introduction}

Large language models are increasingly limited not by model quality but
by inference cost. Even when a target model fits comfortably on a modern
GPU, autoregressive decoding remains sequential, so throughput and
latency can become the bottleneck for interactive use and large-scale
evaluation. This project studies two complementary acceleration
strategies for the same target model, \texttt{google/gemma-3-12b-it}:

\begin{enumerate}
\def\labelenumi{\arabic{enumi}.}
\item
  a standard speculative-decoding module implemented in
  \texttt{Code/gemma-draft-pair/}, using a smaller draft model to
  propose tokens for target-model verification;
\item
  an EAGLE-3 module implemented in
  \texttt{venkatesh-nagarjuna-individual-project/Code/}, using a learned
  draft head over target hidden states instead of a second full language
  model.
\end{enumerate}

The purpose of the project is not to change the target model itself, but
to improve inference efficiency while preserving practical reliability.
The standard speculative-decoding module evaluates the classical
two-model draft-and-verify design introduced by Leviathan et al. (2023).
The EAGLE-3 module evaluates a more compact alternative inspired by Li
et al. (2024), where future-token proposals are produced from internal
target-model features rather than from a separate draft model.

This report is intentionally restricted to experiments that were
actually completed and saved in the repository. In particular:

\begin{enumerate}
\def\labelenumi{\arabic{enumi}.}
\item
  the fully finished standard speculative-decoding benchmark set is the
  pair \texttt{F} Gemma-3 run in
  \texttt{Code/gemma-draft-pair/gemma\_runs/outputs/F\_final/summary.csv};
\item
  the fully finished EAGLE-3 training and inference artifacts are stored
  under \texttt{venkatesh-nagarjuna-individual-project/Code/artifacts/}.
\end{enumerate}

The rest of the report follows the course guidelines closely. Section 2
describes the datasets. Section 3 explains the models and algorithms.
Section 4 describes the experimental setup and evaluation criteria.
Section 5 discusses searched hyperparameters, overfitting control, and
extrapolation control. Section 6 presents the main results with figures
and tables. Section 7 summarizes the conclusions and practical lessons.
Section 8 lists papers, websites, and repositories used for background
information or implementation support.

\subsection{2. Description of the Data
Set}\label{2-description-of-the-data-set}

\subsubsection{2.1 Evaluation Datasets}\label{21-evaluation-datasets}

Both acceleration modules were evaluated on the same four benchmark
prompt collections. Each configuration used \texttt{50} prompts sampled
with seed \texttt{42}.

{\def\LTcaptype{none} % do not increment counter
\begin{longtable}[]{@{}llllll@{}}
\toprule\noalign{}
Task & Hugging Face dataset & Split & Prompt field & Purpose & Max
prompt tokens \\
\midrule\noalign{}
\endhead
\bottomrule\noalign{}
\endlastfoot
\texttt{humaneval} & \texttt{openai/openai\_humaneval} & \texttt{test} &
\texttt{prompt} & Code completion & 512 \\
\texttt{triviaqa} & \texttt{mandarjoshi/trivia\_qa} (\texttt{rc}) &
\texttt{validation} & \texttt{question} & Short-form question answering
& 256 \\
\texttt{cnn\_dailymail} & \texttt{abisee/cnn\_dailymail}
(\texttt{3.0.0}) & \texttt{test} & \texttt{article} & Summarization &
1024 \\
\texttt{writingprompts} & \texttt{euclaise/writingprompts} &
\texttt{validation} & \texttt{prompt} & Creative writing continuation &
256 \\
\end{longtable}
}

The prompt-processing pipeline applies the tokenizer\textquotesingle s
chat template when one is available and falls back to a simple
\texttt{System\ /\ User\ /\ Assistant} text format otherwise. This kept
evaluation consistent across model families and reduced
prompt-formatting differences as a confound.

These four tasks were intentionally heterogeneous. \texttt{humaneval}
rewards deterministic, syntax-sensitive continuation; \texttt{triviaqa}
focuses on short-answer generation; \texttt{cnn\_dailymail} stresses
long-context summarization; and \texttt{writingprompts} tests
longer-form open-ended generation. This mix is useful because
acceleration quality depends strongly on how predictable the next-token
distribution is for a given task.

\subsubsection{2.2 EAGLE-3 Training
Dataset}\label{22-eagle-3-training-dataset}

The EAGLE-3 draft head was trained on \texttt{vicgalle/alpaca-gpt4}.
Each example was converted into an instruction-following chat sequence
with a fixed system prompt, a user message formed from the instruction
and optional input, and the reference assistant output.

{\def\LTcaptype{none} % do not increment counter
\begin{longtable}[]{@{}lllll@{}}
\toprule\noalign{}
Dataset & Split & Fields used & Max sequence length & Observed training
size \\
\midrule\noalign{}
\endhead
\bottomrule\noalign{}
\endlastfoot
\texttt{vicgalle/alpaca-gpt4} & \texttt{train} & \texttt{instruction},
\texttt{input}, \texttt{output} & 512 & 52,002 examples \\
\end{longtable}
}

The completed training log confirms that \texttt{52,002} examples were
prepared. Importantly, the EAGLE-3 training data is distinct from the
four held-out evaluation tasks above, so the final benchmark results are
not simple memorization of the supervised training set.

\subsection{3. Description of the NLP Model and
Algorithm}\label{3-description-of-the-nlp-model-and-algorithm}

\subsubsection{3.1 Target Model}\label{31-target-model}

Both modules use \texttt{google/gemma-3-12b-it} as the target model. In
both cases, the target model remains the generator whose behavior we
ultimately care about. The acceleration logic changes how candidate
tokens are proposed and verified, but the target model remains
responsible for the final sequence.

The project uses 4-bit target loading for practical GPU memory reasons.
This choice does not change the task definitions or the decoding
algorithms, but it makes the Gemma-3-12B target usable on an AWS
\texttt{g5.2xlarge} instance with an NVIDIA A10G GPU.

The broader codebase also contains a Gemma-4-31B comparison path, even
though the fully completed benchmark results in this report are centered
on Gemma-3-12B. The two targets are therefore relevant at the
model-context level even though they do not contribute equally to the
finished experimental record.

\pandocbounded{\includegraphics[keepaspectratio,alt={}]{C:/Users/venka/OneDrive/Desktop/Class/6312/Final Project/Final-Group-Project-Report/assets/gemma3_gemma4_architecture_overview.png}}

Figure 1 summarizes the model-family context in which the acceleration
methods operate. In this project, Gemma-3-12B is the main target used
for the completed standard speculative-decoding and EAGLE-3 experiments.
Gemma-4-31B is still relevant because the repository supports it as a
larger comparison target, but it was not the center of the finished
benchmark analysis reported here.

\subsubsection{3.2 Standard Speculative
Decoding}\label{32-standard-speculative-decoding}

Standard speculative decoding uses two distributions:

\begin{enumerate}
\def\labelenumi{\arabic{enumi}.}
\item
  the target distribution
  \texttt{p(x\_t\ \textbackslash{}mid\ x\_\{\textless{}t\})}, produced
  by the large model;
\item
  the draft distribution
  \texttt{q(x\_t\ \textbackslash{}mid\ x\_\{\textless{}t\})}, produced
  by the smaller model.
\end{enumerate}

The draft model proposes \texttt{\textbackslash{}gamma} future tokens,
and the target model verifies them in a batched pass. A drafted token
\texttt{x\_t} is accepted with probability:

\[\alpha_t = \min\left(1, \frac{p(x_t \mid x_{<t})}{q(x_t \mid x_{<t})}\right)\]

If a proposal is rejected, the implementation samples from the residual
distribution:

\[r(x) = \frac{\max(0, p(x) - q(x))}{\sum_{x'} \max(0, p(x') - q(x'))}\]

At \texttt{temperature\ =\ 0.0}, the code uses greedy verification
instead of stochastic rejection sampling. The main benefit of the method
is that the target model can verify multiple proposed positions in one
forward pass. The main risk is that speedup collapses when the draft
model is poorly aligned with the target and many proposals are rejected.

The completed configuration reported in this project is pair \texttt{F}.

{\def\LTcaptype{none} % do not increment counter
\begin{longtable}[]{@{}lllll@{}}
\toprule\noalign{}
Pair & Target model & Draft model & Quantization & Estimated VRAM \\
\midrule\noalign{}
\endhead
\bottomrule\noalign{}
\endlastfoot
\texttt{F} & \texttt{google/gemma-3-12b-it} &
\texttt{google/gemma-3-1b-it} & 4-bit target, BF16 draft & about 8.6
GB \\
\end{longtable}
}

This module required no training. Its algorithmic development came from
adapting the speculative-decoding literature to Gemma-specific
engineering constraints such as KV-cache handling, preallocated token
buffers, 4-bit target loading, and detailed timing and VRAM
instrumentation.

\pandocbounded{\includegraphics[keepaspectratio,alt={}]{C:/Users/venka/OneDrive/Desktop/Class/6312/Final Project/Final-Group-Project-Report/assets/speculative_decoding_overview.png}}

Figure 2 illustrates the core intuition behind standard speculative
decoding: if the draft model is often correct, then several future
tokens can be proposed cheaply and verified in fewer expensive
target-model passes. The completed pair \texttt{F} experiments test
exactly how often that assumption holds for Gemma-3-12B paired with
Gemma-3-1B.

\subsubsection{3.3 EAGLE-3 Draft-Head
Acceleration}\label{33-eagle-3-draft-head-acceleration}

EAGLE-3 removes the separate draft model and replaces it with a learned
draft head operating on the target model\textquotesingle s hidden
states. In the completed implementation, the draft head uses:

\begin{enumerate}
\def\labelenumi{\arabic{enumi}.}
\item
  hidden states from three target layers;
\item
  a fusion projection over those three feature streams;
\item
  the target embedding layer for the currently verified token;
\item
  an input projection that combines token and feature information;
\item
  one copied decoder layer;
\item
  the frozen target normalization and LM head.
\end{enumerate}

If the hidden states selected from the low-, mid-, and high-level layers
are denoted by \texttt{h\_low}, \texttt{h\_mid}, and \texttt{h\_high},
then the fused representation is

\[h_{\text{fused}} = W_f [h_{\text{low}}; h_{\text{mid}}; h_{\text{high}}]\]

and the decoder-layer input is

\[z_t = W_{\text{in}} [e(x_t); h_{\text{fused}, t}]\]

where \texttt{e(x\_t)} is the embedding of the current token.

The training objective is a multi-step KL-divergence loss:

\[\mathcal{L} = \sum_{k=0}^{K-1} \lambda^k \, KL\left(p^{(k)}_{\text{target}} \parallel p^{(k)}_{\text{draft}}\right)\]

with default multi-step horizon \texttt{K\ =\ 5} and decay
\texttt{\textbackslash{}lambda\ =\ 0.8}.

This objective is important because EAGLE-3 is not only trying to match
the target\textquotesingle s next-token distribution at one step. It is
also trying to remain stable after its own predictions are fed back in
future steps. In other words, the multi-step loss is explicitly designed
to reduce compounding error during extrapolation.

The completed production configuration is pair \texttt{I}.

{\def\LTcaptype{none} % do not increment counter
\begin{longtable}[]{@{}llll@{}}
\toprule\noalign{}
Pair & Target model & Draft component & Quantization \\
\midrule\noalign{}
\endhead
\bottomrule\noalign{}
\endlastfoot
\texttt{I} & \texttt{google/gemma-3-12b-it} & trained EAGLE-3 draft head
& 4-bit target \\
\end{longtable}
}

The decode path used in this repository is a best-path EAGLE-3 variant
rather than a full masked-tree verification implementation. The code
samples a root token from the target, expands a draft tree from fused
hidden features, extracts candidate paths, chooses the best path, and
verifies that path sequentially with the target model.

\pandocbounded{\includegraphics[keepaspectratio,alt={}]{C:/Users/venka/OneDrive/Desktop/Class/6312/Final Project/Final-Group-Project-Report/assets/eagle_fig1.png}}

Figure 3 is not a literal diagram of the exact code path in this
repository. Instead, it provides the high-level conceptual background
for EAGLE-style acceleration: future-token proposals are generated from
internal target-model representations rather than from a second full
language model. The project's EAGLE-3 implementation specializes this
idea to a Gemma-3-12B target with a learned fused-feature draft head.

\subsection{4. Experimental Setup}\label{4-experimental-setup}

\subsubsection{4.1 Hardware and Software
Framework}\label{41-hardware-and-software-framework}

The project was implemented in Python using PyTorch and Hugging Face
Transformers, with Hugging Face Datasets for benchmark loading and
bitsandbytes for 4-bit target quantization and 8-bit optimizer support.
Plotting used Matplotlib and Seaborn.

The EAGLE-3 training and evaluation logs explicitly record an NVIDIA
A10G GPU, consistent with the AWS \texttt{g5.2xlarge} hardware target
used throughout the project. This is a realistic setting for
model-serving experiments because it forces memory-efficient design
decisions rather than relying on a very large accelerator.

\subsubsection{4.2 Standard Speculative-Decoding
Protocol}\label{42-standard-speculative-decoding-protocol}

The completed standard speculative benchmark for pair \texttt{F} used:

\begin{enumerate}
\def\labelenumi{\arabic{enumi}.}
\item
  \texttt{\textbackslash{}gamma\ \textbackslash{}in\ \textbackslash{}\{1,\ 3,\ 5,\ 7,\ 10\textbackslash{}\}};
\item
  \texttt{temperature\ \textbackslash{}in\ \textbackslash{}\{0.0,\ 0.6,\ 1.0\textbackslash{}\}};
\item
  \texttt{4} tasks;
\item
  \texttt{50} prompts per configuration;
\item
  \texttt{128} maximum generated tokens;
\item
  \texttt{3} warm-up generations before measurement.
\end{enumerate}

This yields \texttt{60} speculative configurations and \texttt{12}
matching baselines, for \texttt{72} saved rows in the completed summary
file.

\subsubsection{4.3 EAGLE-3 Training
Protocol}\label{43-eagle-3-training-protocol}

The completed EAGLE-3 training run used:

\begin{enumerate}
\def\labelenumi{\arabic{enumi}.}
\item
  \texttt{google/gemma-3-12b-it} as a frozen 4-bit target;
\item
  the Alpaca-GPT4 instruction dataset described in Section 2.2;
\item
  BF16 mixed precision for the trainable head;
\item
  8-bit AdamW when available;
\item
  gradient accumulation and gradient clipping;
\item
  periodic checkpointing every \texttt{500} steps.
\end{enumerate}

The finished log shows that the run resumed from \texttt{step5000} and
completed at \texttt{step19503}. The active GPU training time across the
two logged sessions was \texttt{36.98} hours.

\subsubsection{4.4 EAGLE-3 Evaluation
Protocol}\label{44-eagle-3-evaluation-protocol}

The completed EAGLE-3 inference sweep used:

\begin{enumerate}
\def\labelenumi{\arabic{enumi}.}
\item
  \texttt{tree\_budget\ \textbackslash{}in\ \textbackslash{}\{20,\ 60\textbackslash{}\}};
\item
  \texttt{temperature\ \textbackslash{}in\ \textbackslash{}\{0.0,\ 0.6,\ 1.0\textbackslash{}\}};
\item
  the same \texttt{4} evaluation tasks;
\item
  \texttt{50} prompts per configuration;
\item
  \texttt{128} maximum generated tokens;
\item
  \texttt{3} warm-up generations before measurement.
\end{enumerate}

This yields \texttt{24} EAGLE configurations and \texttt{12} matching
baselines, for \texttt{36} rows in the finished summary file.

The long-context \texttt{cnn\_dailymail} cases were particularly
important because they stress KV-cache management. The targeted rerun
log for the \texttt{cnn\_dailymail} fix shows that the implementation
can fall back to uncached verification when cache-shape mismatches
occur. This fallback preserves correctness but reduces speed, which
becomes relevant in the final results.

\subsubsection{4.5 Performance Criteria and Correctness
Checks}\label{45-performance-criteria-and-correctness-checks}

The project evaluated performance primarily as a systems problem. The
two main metrics were tokens per second and speedup over the matching
baseline:

\[TPS = \frac{\text{total generated tokens}}{\text{wall-clock time}}\]

\[Speedup = \frac{TPS_{\text{accelerated}}}{TPS_{\text{baseline}}}\]

The code also logged:

\begin{enumerate}
\def\labelenumi{\arabic{enumi}.}
\item
  time to first token (TTFT),
\item
  mean acceptance rate,
\item
  mean acceptance length,
\item
  draft-overhead ratio,
\item
  peak VRAM.
\end{enumerate}

In addition to these metrics, the codebase includes correctness and
smoke-test scripts. The standard speculative module includes unit tests
and greedy-equivalence tests. The EAGLE module includes unit tests for
tree logic and smoke tests for the active pair \texttt{I}. The final
report emphasizes the completed benchmark outputs, but these tests were
useful for ensuring the accelerators were not simply fast because they
were broken.

\subsection{5. Hyperparameter Search, Overfitting, and Extrapolation
Control}\label{5-hyperparameter-search-overfitting-and-extrapolation-control}

\subsubsection{5.1 Searched
Hyperparameters}\label{51-searched-hyperparameters}

The searched hyperparameters were different for the two modules.

{\def\LTcaptype{none} % do not increment counter
\begin{longtable}[]{@{}lll@{}}
\toprule\noalign{}
Module & Searched hyperparameters & Values \\
\midrule\noalign{}
\endhead
\bottomrule\noalign{}
\endlastfoot
Standard speculative decoding & Speculation length
\texttt{\textbackslash{}gamma} & \texttt{1,\ 3,\ 5,\ 7,\ 10} \\
Standard speculative decoding & Temperature &
\texttt{0.0,\ 0.6,\ 1.0} \\
EAGLE-3 inference & Tree budget & \texttt{20,\ 60} \\
EAGLE-3 inference & Temperature & \texttt{0.0,\ 0.6,\ 1.0} \\
\end{longtable}
}

The EAGLE-3 training run did not perform a large hyperparameter sweep
because training even a lightweight draft head on Gemma-3-12B is
expensive. Instead, one practical configuration was chosen and monitored
carefully.

{\def\LTcaptype{none} % do not increment counter
\begin{longtable}[]{@{}ll@{}}
\toprule\noalign{}
EAGLE-3 training setting & Active value \\
\midrule\noalign{}
\endhead
\bottomrule\noalign{}
\endlastfoot
Learning rate & \texttt{3e-4} \\
Batch size & \texttt{1} \\
Gradient accumulation & \texttt{8} \\
Weight decay & \texttt{0.01} \\
Max sequence length & \texttt{512} \\
Multi-step horizon \texttt{K} & \texttt{5} \\
Multi-step decay \texttt{\textbackslash{}lambda} & \texttt{0.8} \\
\end{longtable}
}

\subsubsection{5.2 Overfitting Prevention and
Detection}\label{52-overfitting-prevention-and-detection}

Overfitting affects the two modules differently.

For standard speculative decoding, there is no trainable component, so
overfitting in the usual supervised-learning sense does not apply. The
main concern is draft-target mismatch rather than training-set
memorization.

For EAGLE-3 training, the project used several practical controls:

\begin{enumerate}
\def\labelenumi{\arabic{enumi}.}
\item
  weight decay and gradient clipping to reduce unstable parameter
  growth;
\item
  warmup and cosine decay to stabilize optimization;
\item
  periodic checkpoints so runs could be resumed rather than restarted;
\item
  evaluation on four held-out benchmark datasets that differ
  substantially from Alpaca-GPT4.
\end{enumerate}

The project did not run a separate validation-perplexity study for the
draft head. That is a real limitation. Overfitting detection therefore
relied on two signals: the smooth downward training-loss curve and the
fact that downstream inference was evaluated on held-out tasks from a
different data distribution.

\subsubsection{5.3 Extrapolation and Compounding-Error
Control}\label{53-extrapolation-and-compounding-error-control}

Extrapolation is especially important for acceleration methods because
both the draft model and the EAGLE draft head are asked to predict
future tokens before the target model sees them.

The project addressed this in different ways:

\begin{enumerate}
\def\labelenumi{\arabic{enumi}.}
\item
  Standard speculative decoding uses target-model verification at every
  round. Poor draft proposals can be rejected immediately.
\item
  EAGLE-3 uses a multi-step KL objective so the head learns to recover
  after its own predictions are fed back in future steps.
\item
  The EAGLE implementation includes a long-context fallback path that
  switches to uncached verification when KV-cache trimming becomes
  unreliable.
\end{enumerate}

This last point matters in practice. The \texttt{cnn\_dailymail} rerun
log shows frequent uncached-fallback warnings at long contexts. The
fallback preserves functional correctness, but it also explains part of
the throughput loss on summarization tasks.

\subsection{6. Results}\label{6-results}

\subsubsection{6.1 Standard Speculative-Decoding
Results}\label{61-standard-speculative-decoding-results}

The completed pair \texttt{F} benchmark is the strongest finished
standard speculative artifact in the repository and serves as the main
evidence for Module 1.

Table 1 summarizes the best speculative configuration for each task
together with the matching baseline.

{\def\LTcaptype{none} % do not increment counter
\begin{longtable}[]{@{}lllllllll@{}}
\toprule\noalign{}
Task & Matching baseline TPS & Best speculative config & Best
speculative TPS & Speedup & Mean acceptance rate & Baseline TTFT (ms) &
Best TTFT (ms) & Peak VRAM (GB) \\
\midrule\noalign{}
\endhead
\bottomrule\noalign{}
\endlastfoot
\texttt{humaneval} & 9.63 & \texttt{\textbackslash{}gamma\ =\ 10},
\texttt{T\ =\ 0.0} & 14.22 & 1.476 & 0.9279 & 179.12 & 555.24 & 9.40 \\
\texttt{triviaqa} & 9.63 & \texttt{\textbackslash{}gamma\ =\ 5},
\texttt{T\ =\ 0.0} & 12.47 & 1.295 & 0.8751 & 134.99 & 320.48 & 9.31 \\
\texttt{cnn\_dailymail} & 9.34 & \texttt{\textbackslash{}gamma\ =\ 1},
\texttt{T\ =\ 0.0} & 8.30 & 0.889 & 0.5342 & 462.10 & 525.01 & 9.77 \\
\texttt{writingprompts} & 9.68 & \texttt{\textbackslash{}gamma\ =\ 1},
\texttt{T\ =\ 0.0} & 8.43 & 0.871 & 0.4525 & 142.32 & 180.03 & 9.32 \\
\end{longtable}
}

The main conclusion from Table 1 is that pair \texttt{F} is useful but
not universally useful. It produces strong gains on \texttt{humaneval}
and \texttt{triviaqa}, but not on \texttt{cnn\_dailymail} or
\texttt{writingprompts}.

\pandocbounded{\includegraphics[keepaspectratio,alt={}]{C:/Users/venka/OneDrive/Desktop/Class/6312/Final Project/Final-Group-Project-Report/assets/speculative_pair_f_speedup_vs_gamma.png}}

Figure 4 shows that speedup improves with larger
\texttt{\textbackslash{}gamma} only when draft-target agreement remains
strong. For \texttt{humaneval}, speedup increases steadily up to
\texttt{1.476x} at \texttt{\textbackslash{}gamma\ =\ 10}. For
\texttt{triviaqa}, the optimum occurs earlier at
\texttt{\textbackslash{}gamma\ =\ 5}. In contrast,
\texttt{cnn\_dailymail} and \texttt{writingprompts} degrade as
\texttt{\textbackslash{}gamma} increases because more aggressive
speculative proposals produce too much verification waste.

\pandocbounded{\includegraphics[keepaspectratio,alt={}]{C:/Users/venka/OneDrive/Desktop/Class/6312/Final Project/Final-Group-Project-Report/assets/speculative_pair_f_acceptance_heatmap.png}}

Figure 5 explains the pattern in Figure 4. Acceptance is consistently
high on \texttt{humaneval} and \texttt{triviaqa}, especially at larger
\texttt{\textbackslash{}gamma}, but substantially lower on
\texttt{cnn\_dailymail} and \texttt{writingprompts}. The strongest
speedups occur where acceptance is both high and stable. Low acceptance
does not simply reduce marginal gains; it can erase them entirely.

\pandocbounded{\includegraphics[keepaspectratio,alt={}]{C:/Users/venka/OneDrive/Desktop/Class/6312/Final Project/Final-Group-Project-Report/assets/speculative_pair_f_best_vs_baseline.png}}

Figure 6 shows the central trade-off of standard speculative decoding.
The same settings that improve throughput also increase first-token
latency. The gain on \texttt{humaneval} and \texttt{triviaqa} is real
and substantial, but it comes with noticeably worse TTFT. This means
pair \texttt{F} is more attractive for long-form or sustained generation
than for applications where first-token responsiveness is the primary
objective.

Across the full finished sweep, peak VRAM remained below
\texttt{10\ GB}, so pair \texttt{F} is operationally lightweight on an
A10G. This is an important positive result: even though speculative
decoding is not beneficial on every task, the cost of trying it on
Gemma-3-12B is modest from a memory perspective.

\subsubsection{6.2 EAGLE-3 Training
Results}\label{62-eagle-3-training-results}

The EAGLE-3 training artifacts provide a clear record of successful
optimization rather than an incomplete training attempt.

Table 2 summarizes the completed run.

{\def\LTcaptype{none} % do not increment counter
\begin{longtable}[]{@{}ll@{}}
\toprule\noalign{}
Item & Observed value \\
\midrule\noalign{}
\endhead
\bottomrule\noalign{}
\endlastfoot
Target model & \texttt{google/gemma-3-12b-it} \\
Logged hardware & \texttt{NVIDIA\ A10G} \\
Prepared training examples & \texttt{52,002} \\
Trainable parameters & \texttt{297,876,992} \\
Optimizer & 8-bit AdamW \\
Final checkpoint step & \texttt{19,503} \\
Resume point & \texttt{step5000} \\
Active GPU training time & \texttt{36.98} hours \\
Final checkpoint & \texttt{eagle3\_gemma3\_12b\_final.pt} \\
\end{longtable}
}

Selected loss milestones are shown in Table 3.

{\def\LTcaptype{none} % do not increment counter
\begin{longtable}[]{@{}ll@{}}
\toprule\noalign{}
Milestone & Loss \\
\midrule\noalign{}
\endhead
\bottomrule\noalign{}
\endlastfoot
Step 10 & 145.2908 \\
Step 100 & 21.2564 \\
Step 500 & 9.7899 \\
Epoch 1 average loss & 1.8044 \\
Epoch 2 average loss & 1.2272 \\
Step 19,000 & 0.7941 \\
Epoch 3 average loss & 0.8095 \\
\end{longtable}
}

\pandocbounded{\includegraphics[keepaspectratio,alt={}]{C:/Users/venka/OneDrive/Desktop/Class/6312/Final Project/Final-Group-Project-Report/assets/eagle3_training_loss.png}}

Figure 7 shows a very steep initial loss drop followed by a long
convergence tail. This is the behavior expected from a large supervised
approximation problem: the head quickly learns coarse alignment with the
target distribution and then spends most of the remaining optimization
budget improving fine-grained agreement. The end of the run remains well
below \texttt{1.0} average loss, which is strong evidence that the head
became a useful proposal mechanism rather than a random extrapolator.

The training result is important because it establishes that the project
did not stop at benchmarking a training-free baseline. The EAGLE-3
module was trained end to end, resumed successfully after an
interruption, and finished with a saved final checkpoint.

\subsubsection{6.3 EAGLE-3 Inference
Results}\label{63-eagle-3-inference-results}

The completed EAGLE-3 inference sweep provides a full task-by-task
comparison against matching baselines. Table 4 summarizes the best EAGLE
configuration for each task.

{\def\LTcaptype{none} % do not increment counter
\begin{longtable}[]{@{}llllllllll@{}}
\toprule\noalign{}
Task & Best tree budget & Best temperature & Matching baseline TPS &
Best EAGLE TPS & Speedup & Mean acceptance rate & Matching baseline TTFT
(ms) & Best EAGLE TTFT (ms) & Best EAGLE peak VRAM (GB) \\
\midrule\noalign{}
\endhead
\bottomrule\noalign{}
\endlastfoot
\texttt{humaneval} & 20 & 1.0 & 8.52 & 7.61 & 0.893 & 0.0974 & 203.77 &
202.53 & 8.29 \\
\texttt{triviaqa} & 20 & 0.0 & 8.82 & 20.87 & 2.368 & 0.8977 & 162.03 &
162.60 & 8.11 \\
\texttt{cnn\_dailymail} & 20 & 0.6 & 8.30 & 6.03 & 0.727 & 0.0803 &
475.55 & 476.55 & 9.25 \\
\texttt{writingprompts} & 20 & 0.0 & 8.27 & 8.01 & 0.969 & 0.2206 &
172.87 & 169.72 & 8.39 \\
\end{longtable}
}

Table 4 shows a much more uneven task profile than the training-loss
curve alone would suggest. \texttt{triviaqa} is an excellent fit for the
trained EAGLE-3 head, but \texttt{humaneval}, \texttt{cnn\_dailymail},
and \texttt{writingprompts} remain below or near baseline.

\pandocbounded{\includegraphics[keepaspectratio,alt={}]{C:/Users/venka/OneDrive/Desktop/Class/6312/Final Project/Final-Group-Project-Report/assets/eagle3_final_speedup_heatmap.png}}

Figure 8 makes the task dependence of EAGLE-3 especially clear.
\texttt{triviaqa} is the standout case, peaking at \texttt{2.368x}
speedup for \texttt{tree\_budget\ =\ 20} and
\texttt{temperature\ =\ 0.0}. \texttt{writingprompts} comes close to
parity, \texttt{humaneval} improves somewhat at higher temperatures but
remains below \texttt{1.0x}, and \texttt{cnn\_dailymail} remains well
below baseline throughout the sweep.

An additional practical finding from Figure 8 is that
\texttt{tree\_budget\ =\ 20} outperforms \texttt{tree\_budget\ =\ 60}
for every task\textquotesingle s best finished result. Larger trees did
not help once verification cost was included.

\pandocbounded{\includegraphics[keepaspectratio,alt={}]{C:/Users/venka/OneDrive/Desktop/Class/6312/Final Project/Final-Group-Project-Report/assets/eagle3_final_acceptance_heatmap.png}}

Figure 9 explains why \texttt{triviaqa} is the only task with a large
EAGLE speedup. Its acceptance rate is far higher than the other tasks,
especially at low temperature. \texttt{humaneval} is the opposite
extreme: it achieves very low acceptance even when its TTFT remains
competitive with baseline. \texttt{cnn\_dailymail} and
\texttt{writingprompts} sit in the middle but do not maintain enough
agreement to overcome verification overhead.

\pandocbounded{\includegraphics[keepaspectratio,alt={}]{C:/Users/venka/OneDrive/Desktop/Class/6312/Final Project/Final-Group-Project-Report/assets/eagle3_best_vs_baseline_tps.png}}

Figure 10 provides the clearest direct comparison between EAGLE and
baseline throughput. The \texttt{triviaqa} improvement is not marginal;
it is the dominant success case of the entire EAGLE study. At the same
time, the figure also shows that the method should not be treated as
uniformly beneficial. The same draft head that works extremely well on
\texttt{triviaqa} is a net slowdown on summarization.

\pandocbounded{\includegraphics[keepaspectratio,alt={}]{C:/Users/venka/OneDrive/Desktop/Class/6312/Final Project/Final-Group-Project-Report/assets/eagle3_best_ttft_vram.png}}

Figure 11 shows that the best EAGLE configurations remain practical from
a deployment perspective. TTFT stays in roughly the same scale as
baseline for all tasks, and peak VRAM remains in the \texttt{8.11} to
\texttt{9.25\ GB} range. The slowdown on \texttt{cnn\_dailymail} is
therefore not a memory-capacity issue. It is primarily an
acceptance-and-verification issue, made worse by long-context fallback
behavior.

The targeted \texttt{cnn\_dailymail} rerun log supports this
interpretation. It shows repeated uncached-verification warnings at long
contexts, followed by successful completion. In other words, the current
implementation prioritizes correctness and robustness over raw speed
when the cache path becomes unreliable. This is the right engineering
choice, but it also explains why summarization remains the hardest case.

\subsection{7. Summary and Conclusions}\label{7-summary-and-conclusions}

This project produced two useful and complementary conclusions about
inference acceleration for Gemma-3-12B.

First, standard speculative decoding with a separate small draft model
is already effective on some tasks. Pair \texttt{F} achieved the
strongest finished result in the repository on \texttt{humaneval}, where
the best configuration reached \texttt{1.476x} speedup, and it also
improved \texttt{triviaqa} to \texttt{1.295x}. These gains were strongly
correlated with high acceptance rates. However, the same method slowed
down \texttt{cnn\_dailymail} and \texttt{writingprompts}, demonstrating
that speculative decoding is only worthwhile when the draft model tracks
the target closely enough.

Second, the EAGLE-3 draft-head approach is viable and can outperform
standard drafting on selected tasks while using only the target model
and a lightweight learned head. The completed training run converged
successfully, produced a final checkpoint, and yielded a striking
\texttt{2.368x} speedup on \texttt{triviaqa}. That is the best
acceleration result in the entire project. At the same time, the EAGLE
results also show that a trained hidden-state extrapolator is not
universally strong across domains. Long-context summarization remains
difficult, and low-acceptance tasks still pay too much verification
overhead.

The most important lessons from the project are:

\begin{enumerate}
\def\labelenumi{\arabic{enumi}.}
\item
  draft-target agreement is the dominant driver of speedup;
\item
  throughput gains and first-token latency do not necessarily improve
  together;
\item
  a lightweight learned draft head can be more effective than a separate
  draft model on the right task;
\item
  cache management and long-context robustness matter as much as the
  nominal decoding algorithm;
\item
  realistic GPU-memory limits force design choices that are central
  rather than incidental.
\end{enumerate}

The most natural future improvements are:

\begin{enumerate}
\def\labelenumi{\arabic{enumi}.}
\item
  adding a validation-driven EAGLE training study rather than relying on
  one production training configuration;
\item
  improving long-context verification so \texttt{cnn\_dailymail} does
  not rely on slow fallback behavior;
\item
  investigating whether task-adaptive tree budgets or confidence
  thresholds can improve EAGLE robustness;
\item
  expanding pair-\texttt{F} style speculative experiments with more
  systematically matched draft models.
\end{enumerate}

Overall, the project shows that Gemma-3-12B inference can be accelerated
substantially, but only when the proposal mechanism matches the task
structure. There is no single universally best accelerator. Instead, the
finished experiments support a more nuanced conclusion: standard
speculative decoding is a practical low-risk baseline, while EAGLE-3
offers the most exciting upside when the hidden-state extrapolator is
well aligned with the target distribution.

\subsection{8. References}\label{8-references}

\begin{enumerate}
\def\labelenumi{\arabic{enumi}.}
\item
  Leviathan, Y., Kalman, M., and Matias, Y. \emph{Fast Inference from
  Transformers via Speculative Decoding.} ICML 2023.
  \url{https://arxiv.org/abs/2302.01318}
\item
  Li, Y., Wei, F., Zhang, C., and Zhang, H. \emph{EAGLE: Speculative
  Sampling Requires Rethinking Feature Uncertainty.} ICML 2024.
  \url{https://arxiv.org/abs/2401.15077}
\item
  Gemma Team. \emph{Gemma 3 Technical Report.} Google DeepMind, 2025.
  \url{https://storage.googleapis.com/deepmind-media/gemma/Gemma3Report.pdf}
\item
  SafeAILab EAGLE repository. \url{https://github.com/SafeAILab/EAGLE}
\item
  PyTorch repository. \url{https://github.com/pytorch/pytorch}
\item
  Hugging Face Transformers repository.
  \url{https://github.com/huggingface/transformers}
\item
  bitsandbytes repository.
  \url{https://github.com/bitsandbytes-foundation/bitsandbytes}
\item
  Streamlit repository. \url{https://github.com/streamlit/streamlit}
\item
  Hugging Face model card: \texttt{google/gemma-3-1b-it}.
  \url{https://huggingface.co/google/gemma-3-1b-it}
\item
  Hugging Face model card: \texttt{google/gemma-3-12b-it}.
  \url{https://huggingface.co/google/gemma-3-12b-it}
\item
  Hugging Face dataset: \texttt{openai/openai\_humaneval}.
  \url{https://huggingface.co/datasets/openai/openai_humaneval}
\item
  Hugging Face dataset: \texttt{mandarjoshi/trivia\_qa}.
  \url{https://huggingface.co/datasets/mandarjoshi/trivia_qa}
\item
  Hugging Face dataset: \texttt{abisee/cnn\_dailymail}.
  \url{https://huggingface.co/datasets/abisee/cnn_dailymail}
\item
  Hugging Face dataset: \texttt{euclaise/writingprompts}.
  \url{https://huggingface.co/datasets/euclaise/writingprompts}
\item
  Hugging Face dataset: \texttt{vicgalle/alpaca-gpt4}.
  \url{https://huggingface.co/datasets/vicgalle/alpaca-gpt4}
\end{enumerate}

\end{document}
